\documentclass[aspectratio=169,11pt]{beamer}
\usepackage[spanish]{babel}
\usepackage[T1]{fontenc}
\usepackage[utf8]{inputenc}
\usepackage{lmodern}
\usepackage{ragged2e}
\usepackage{xcolor}

\usetheme{Madrid}
\setbeamertemplate{navigation symbols}{}
\setbeamertemplate{footline}[frame number]

\definecolor{azul}{RGB}{33,87,146}
\setbeamercolor{structure}{fg=azul}
\setbeamercolor{frametitle}{fg=white,bg=azul}
\setbeamercolor{title}{fg=azul}

\title{Semana 04: Dimensiones y Diseno}
\author{Grupo 6}
\date{CIB12 - Aprendizaje de Maquina y Mineria de Datos}

\begin{document}

\begin{frame}
  \titlepage
\end{frame}

\begin{frame}{Tema del Grupo}
\justifying
\textbf{Tema tentativo:}

\vspace{0.5em}

Modelo de aprendizaje profundo con mecanismos de atencion para mejorar la calidad del servicio de guia virtual en servicios de guia virtual con aplicaciones moviles e IoT con senales Wi-Fi/BLE en museos y espacios culturales, 2026.

\vspace{0.7em}

\textbf{Objetivo de la semana 04:}
\begin{itemize}
  \item definir variables;
  \item proponer dimensiones e indicadores;
  \item plantear el diseno preliminar del estudio.
\end{itemize}
\end{frame}

\begin{frame}{Variable Independiente}
\justifying
\textbf{Variable independiente:} Modelo de aprendizaje profundo con mecanismos de atencion.

\vspace{0.6em}

\textbf{Dimensiones:}
\begin{itemize}
  \item arquitectura del modelo;
  \item mecanismo de atencion;
  \item robustez frente a ruido y heterogeneidad;
  \item capacidad de generalizacion.
\end{itemize}

\textbf{Indicadores sugeridos:}
\begin{itemize}
  \item accuracy por zona o sala;
  \item error medio de localizacion;
  \item F1-score;
  \item estabilidad entre escenarios o dispositivos.
\end{itemize}
\end{frame}

\begin{frame}{Variable Dependiente}
\justifying
\textbf{Variable dependiente:} Calidad del servicio de guia virtual.

\vspace{0.6em}

\textbf{Dimensiones:}
\begin{itemize}
  \item precision de localizacion;
  \item oportunidad de activacion;
  \item pertinencia contextual;
  \item continuidad del servicio.
\end{itemize}

\textbf{Indicadores sugeridos:}
\begin{itemize}
  \item porcentaje de activaciones correctas;
  \item tiempo de respuesta;
  \item tasa de error por sala o zona;
  \item porcentaje de sesiones con guia contextual correcta.
\end{itemize}
\end{frame}

\begin{frame}{Desempeño Tecnico y Servicio}
\justifying
\textbf{Idea clave:}

La precision tecnica de localizacion no es el fin en si mismo. Su valor esta en que permite mejorar la calidad del servicio de guia virtual.

\vspace{0.7em}

\begin{itemize}
  \item El error de localizacion mide el rendimiento tecnico.
  \item La activacion correcta de la guia mide el impacto en el servicio.
  \item Ambas dimensiones deben analizarse juntas.
\end{itemize}
\end{frame}

\begin{frame}{Diseño Preliminar}
\justifying
\textbf{Tipo y nivel:}
\begin{itemize}
  \item investigacion aplicada;
  \item nivel explicativo y experimental;
  \item metodo cuantitativo.
\end{itemize}

\vspace{0.6em}

\textbf{Diseño del estudio:}
\begin{itemize}
  \item comparacion entre baseline y modelo propuesto;
  \item uso de fingerprints Wi-Fi/BLE;
  \item evaluacion en escenarios con heterogeneidad de dispositivos y variaciones temporales.
\end{itemize}
\end{frame}

\begin{frame}{Conexion con la Matriz}
\justifying
Las dimensiones e indicadores definidos en esta semana preparan el siguiente entregable:

\vspace{0.6em}

\begin{itemize}
  \item problemas especificos mas medibles;
  \item objetivos mejor delimitados;
  \item hipotesis observables;
  \item variables ya operacionalizadas para la matriz de consistencia.
\end{itemize}
\end{frame}

\begin{frame}{Siguiente Entregable}
\justifying
\textbf{Entrega actual:} dimensiones y diseño.

\vspace{0.7em}

\textbf{Siguiente entregable:} matriz de consistencia.

\vspace{0.7em}

\textbf{Trabajo inmediato del grupo:}
\begin{itemize}
  \item conectar problemas, objetivos e hipotesis;
  \item consolidar variables, dimensiones e indicadores;
  \item pasar del esquema conceptual a una matriz formal.
\end{itemize}
\end{frame}

\end{document}
